\documentclass[12pt, ukrainian]{article}
\usepackage[a4paper]{geometry}
\setlength\parindent{0mm}
\geometry{top=1.1cm,bottom=1.1cm,left=1.1cm,right=1.1cm}
\pagestyle{empty}
\usepackage[T2A]{fontenc}
\usepackage[utf8]{inputenc}
\usepackage[ukrainian]{babel}
\usepackage{graphicx}
\usepackage{caption}
\usepackage{verbatim, listings, clrscode3e, fancyvrb}
\DeclareCaptionFormat{nocaption}{}
\captionsetup{format=nocaption,aboveskip=0pt,belowskip=0pt}
\usepackage[Export]{adjustbox}
\adjustboxset{max size={0.9\linewidth}{0.9\paperheight}}
\def\code#1{\Verb!#1!}
\begin{document}
{{16.05.2024 Контрольна робота, варіант  \No \textbf { 41 }} група К\_\_ ПІБ {\parbox {2.5 cm} \dotfill}\par}
{
%\begin {large}
\begin{minipage}[t]{9.7cm}
У завданнях 1--3 вказати, що буде виведено за виконання. Повідомлення інтерпретатора про помилку позначати ERROR


               
{\begin{center}Завдання {1}\end{center}}
\vspace{-4mm}

\begin{verbatim}
d = [0,1,2,3,4,5,6,7,8,9]
print(d[5::-1])

\end{verbatim}            

\vspace{15mm}
               
{\begin{center}Завдання {2}\end{center}}
\vspace{-4mm}

\begin{verbatim}
try:
    s = {87:4, 47:6, 88:4, 87:5}
    s[33] = 2
    for x in s.items():
        print(x, sep=' ', end=' ')
except:
    print('ERR')

\end{verbatim}            


\end{minipage}
\vline \hspace{8mm}
\begin{minipage}[t]{8.0cm}                
               
{\begin{center}Завдання {3}\end{center}}
\vspace{-4mm}

\begin{verbatim}
_a = 0

class A:
    _a = 1
    
    def __init__(self):
        self._a = 2
        _a = 3
        A._a = 4

obj = A()
obj._a = 5
try:
    print(_a, end=' ')
    print(obj._a, end=' ')
    print(A._a, end=' ')
    print(obj._a, end=' ')
    print(obj._A__a)
except:
    print('ERR')

\end{verbatim}            

\end{minipage}


               
{\begin{center}Завдання {4}\end{center}}
\vspace{-4mm}
Записати ВИРАЗ, значенням якого є множина, що складається з кубів цілих чисел від 18 до 2011 (включно), що не діляться на 3.\par

\vspace{10mm}
               
{\begin{center}Завдання {5}\end{center}}
\vspace{-4mm}
Змінні \kw{next} та \kw{prev} вузла списку позначають наступний та попередній за ним вузол відповідно. Починаючи з голови, у список записано послідовні цілі числа від~$-36$ до~$33$. Нехай змінна \kw{p} позначає вузол, що зберігає значення~$-4$.
Яке число зберігається у вузлі \kw{p.next.prev.next.next.next} ?\par

               
{\begin{center}Завдання {6}\end{center}}
\vspace{-4mm}
Значенням змінної \kw{a} є цілочисловий масив типу $numpy.ndarray$ розмірів $5{\times}4$. На скільки збільшиться сума елементів масиву \kw{a} після виконання
\begin{verbatim}
a[-3] += 1
\end{verbatim}


Якщо код некоректний, то в якості відповіді зазначити ERROR.\par

               
{\begin{center}Завдання {10}\end{center}}
\vspace{-4mm}

Написати функцію $f$, яка для файлу, заданого шляхом, розв’язує задачу:\\ Текстовий файл містить послідовність чисел $a_1$, $a_2$, $\ldots$, $a_t$, записану у форматі одне число на рядок. Знайти третє найбільше значення, що записане у файлі. Кожне значення, що записане у файлі, враховується тільки один раз. Для послідовності -1, -4, -12, -5, -1, 0 таким значенням буде -4.

В усіх випадках розміри файлів такі, що їх вміст повністю завантажений у пам'ять бути не може. Якість коду також оцінюється.



\vspace{10mm}\par

\newpage
\begin{minipage}[t]{9.7cm}
               
{\begin{center}Завдання {7}\end{center}}
\vspace{-4mm}

\begin{center}
	\adjustimage{max size={0.8\linewidth}{0.8\paperheight}}
	{../15BinTree/trees_exp/23.png}
\end{center}
Указати послідовність міток вершин дерева, зображеного на рис., що утвориться в результаті обходу дерева в прямому порядку. Мітки записувати через пробіл.\par Алгоритм починає роботу з кореня (на рис. це верхня вершина).\par

\end{minipage}
\vline \hspace{8mm}
\begin{minipage}[t]{8.0cm}
               
{\begin{center}Завдання {8}\end{center}}
\vspace{-4mm}
Для графа на рисунку з вершини 0 запущено алгоритм пошуку в глибину, що будує кістякове дерево. Списки суміжності вершин переглядаються алгоритмом за зростанням міток вершин.\par
\begin{center}
	\adjustimage{max size={0.9\linewidth}{0.9\paperheight}}
	{../16Graphs/Traversals/107.png}
\end{center}
\par Які з наведених ребер входять до складу отриманого кістякового дерева?\par
1) (1, 6)\par
2) (3, 5)\par
3) (1, 5)\par
4) (2, 5)\par
5) (2, 6)\par
6) (0, 1)\par
{\vskip 15pt} \par

\end{minipage}

               
{\begin{center}Завдання {9}\end{center}}
\vspace{-4mm}
Файл \code{'1.txt'} містить журнал з рядками, в яких через двокрапку записано поряд\-ко\-вий номер запису в журналі (ціле), тип події (ERROR, WARNING, INFO, STATUS), ім'я користувача (складається з латиниці), час події (фор\-мат як у прикладі), код події (ціле).

Білі символи навколо роздільників не допускаються.

Приклад рядка файлу



\begin{verbatim}
13:ERROR:username:2022-05-05 13-26:404
\end{verbatim}


Нижче наведено код, який шукає усі коди помилок (подія ERROR), що зустріча\-ють\-ся в журналі, та зберігає їх у змінній \kw{codes}.



\begin{verbatim}
codes = set()
regex = re.compile(r'XXX')
with open('1.txt', encoding='utf8') as f:
    for s in f:
        res = re.fullmatch(regex, s)
        if not res: continue
        s = ''
        codes.add(YYY)
\end{verbatim}



Який рядок треба записати на місце \kw{XXX} ?
Який рядок треба записати на місце \kw{YYY} ?\par

%\end {large}
}
%end
\end{document}

